%% Chapter 12: Fine-Grained Lower Bounds
%% Status: drafted; 2 figure placeholders.
%%
%% NEW MATERIAL.  There is no deck for this chapter and no pool question yet.
%% It was written on the author's request for a lecture on ETH and SETH, with
%% the 3-colouring and Frechet distance lower bounds as the two worked
%% examples.  See OVERVIEW.md for the list of further topics that were
%% considered and what each would cost.

\chapter{Fine-Grained Lower Bounds}
\label{ch:lower-bounds}

\Cref{ch:exact-i,ch:exact-ii} produced a table of running times: $1.3217^n$ for
$3$-colouring, $\Ostar{}(2^n)$ for Hamiltonian cycle, $\Ostar{}(2^{n/2})$ for
subset sum. Every one of them is an upper bound, and every one of them raises
the same question --- is this the truth, or have we simply not been clever
enough? $\mathrm{P} \ne \mathrm{NP}$ does not answer it. That conjecture says
these problems need superpolynomial time; it says nothing whatever about
whether $3$-colouring needs $1.3^n$ or $1.0001^n$ or $2^{\sqrt n}$.

This chapter is about the conjectures that do answer it, and about what follows
from them. Two hypotheses --- the exponential time hypothesis and its strong
form --- turn out to imply sharp lower bounds for a long list of concrete
problems, exponential and polynomial alike. The technique is the one from
\textrm{NP}-hardness, with the accounting made exact: it is no longer enough
that the reduction be polynomial, it must not waste the resource we are
measuring.

\begin{goals}
  \goalitem{state the exponential time hypothesis and the strong exponential
        time hypothesis, and explain how they differ from
        $\mathrm{P} \ne \mathrm{NP}$}
  \goalitem{state the sparsification lemma and explain what it is needed for}
  \goalitem{explain what a linear reduction is and why an ordinary polynomial
        reduction is not good enough for a fine-grained lower bound}
  \goalitem{show that, assuming ETH, \lang{3-Coloring} has no
        $2^{o(n)}$ algorithm}
  \goalitem{define \lang{Orthogonal Vectors} and prove that SETH implies it has
        no $\bigoh{N^{2-\delta}}$ algorithm}
  \goalitem{define the discrete Fr\'echet distance, and explain how a lower
        bound for it follows from \lang{Orthogonal Vectors}}
  \goalitem{say what these hypotheses would mean if they turned out to be
        false}
\end{goals}

%==============================================================================
\section{What the hypotheses are for}
\label{sec:eth-motivation}
%==============================================================================

An \textrm{NP}-hardness proof is a coarse instrument. It says a problem has no
polynomial algorithm unless $\mathrm{P} = \mathrm{NP}$, and there it stops. For
almost every question an algorithm designer actually has, that is not enough:

\begin{itemize}[leftmargin=*]
  \item \emph{Is $\Ostar{}(2^n)$ for Hamiltonian cycle the end of the road, or
        should we be looking for $1.5^n$?}
  \item \emph{The dynamic program of \cref{ch:treewidth} solves dominating set
        in $\Ostar{}(3^k)$ on a tree decomposition of width $k$. Is the $3$ an
        artefact of the three states, or is it forced?}
  \item \emph{Edit distance, Fr\'echet distance and longest common subsequence
        all have textbook quadratic algorithms and no better one. Is that a
        gap in our knowledge or a fact about the problems?}
\end{itemize}

None of these is a question about polynomial versus exponential. They are
questions about the constant in the exponent, or about the exponent of a
polynomial, and answering them needs a hypothesis with a constant in it.

\begin{intuition}
It is worth being honest about what is going on here. We are not proving lower
bounds; we are proving that certain algorithms cannot exist \emph{unless} a
particular conjecture is false. The value of this depends entirely on whether
the conjecture is believable, and the case for these two is that they have
survived decades of attention and that a great many independent problems turn
out to rest on them. That is the same standing $\mathrm{P} \ne \mathrm{NP}$
has, one level down.
\end{intuition}

%==============================================================================
\section{ETH and SETH}
\label{sec:eth-seth}
%==============================================================================

Both hypotheses are about $k$-SAT, and both are stated in terms of one
quantity.

\begin{definition}\label{def:sk}
For $k \ge 3$ let
\[ s_k \;=\; \inf\bigl\{ \delta \;:\; \text{$k$-SAT can be solved in }
   \bigoh{2^{\delta n}} \text{ time on formulas with $n$ variables} \bigr\}. \]
\end{definition}

The trivial algorithm gives $s_k \le 1$ for every $k$, and $s_2 = 0$ because
$2$-SAT is in $\mathrm{P}$. Everything else is conjecture.

\begin{definition}[Exponential time hypothesis]\label{def:eth}
$\mathrm{ETH}$: $s_3 > 0$. Equivalently, there is no algorithm solving $3$-SAT
in $2^{o(n)}$ time.
\end{definition}

\begin{definition}[Strong exponential time hypothesis]\label{def:seth}
$\mathrm{SETH}$: $\lim_{k \to \infty} s_k = 1$. Equivalently, for every
$\delta < 1$ there is a $k$ such that $k$-SAT cannot be solved in
$\bigoh{2^{\delta n}}$ time.
\end{definition}

Note the shape of \cref{def:seth}: it does not say that $3$-SAT needs $2^n$
time --- it does not, the best known algorithms run in about $1.307^n$ --- but
that the savings vanish as the clause width grows. In particular SETH implies
that CNF-SAT with unbounded clause width has no $\bigoh{2^{\delta n}}$
algorithm for any $\delta < 1$.

\begin{proposition}\label{prop:eth-seth-p-np}
$\mathrm{SETH}$ implies $\mathrm{ETH}$, and $\mathrm{ETH}$ implies
$\mathrm{P} \ne \mathrm{NP}$. Neither converse is known.
\end{proposition}

\begin{proof}
If $s_3 = 0$ then, since \citet{ImpagliazzoPaturi2001} show that the sequence
$(s_k)$ satisfies $s_k \le c \cdot s_3$ for an absolute constant $c$ --- the
complexity of $k$-SAT is controlled by that of $3$-SAT --- we would have
$s_k = 0$ for every $k$, contradicting SETH. And a $2^{o(n)}$ algorithm for
$3$-SAT is in particular not polynomial-time, so ETH is a stronger assumption
than $\mathrm{P} \ne \mathrm{NP}$.
\end{proof}

\begin{remark}[How much are they believed?]\label{rem:believing-eth}
Not equally. ETH is widely believed. SETH is genuinely doubted by some
researchers: it is a statement about a limit, it would be refuted by any single
CNF-SAT algorithm running in $\bigoh{1.99^n}$, and such algorithms are known
for restricted classes of formulas. A refutation of SETH would be a major
result rather than a catastrophe --- it would falsify a long list of
conditional lower bounds without touching any algorithm. This is worth saying
to students, because a conditional lower bound is only as interesting as its
hypothesis, and the two hypotheses in this chapter are not on the same footing.
\end{remark}

%==============================================================================
\section{Reductions that do not waste}
\label{sec:linear-reductions}
%==============================================================================

Here is the difficulty. \lang{3-Coloring} is \textrm{NP}-hard, so there is a
polynomial reduction from $3$-SAT to it. That alone gives \emph{nothing} about
$2^{o(n)}$ algorithms: if the reduction turns an $n$-variable formula into a
graph on $n^2$ vertices, then a $2^{o(n)}$ algorithm for $3$-colouring gives
only $2^{o(n^2)}$ for $3$-SAT, which ETH does not forbid.

So the reduction has to be measured, not merely bounded.

\begin{definition}[Linear reduction]\label{def:linear-reduction}
A reduction from problem $A$ to problem $B$ is \emph{linear} if it runs in
polynomial time and maps an instance whose parameter of interest is $n$ to an
instance whose parameter of interest is $\bigoh{n}$.
\end{definition}

\begin{lemma}\label{lem:linear-transfer}
If there is a linear reduction from $A$ to $B$ and $B$ has a $2^{o(n)}$
algorithm, then so does $A$.
\end{lemma}

\begin{proof}
Run the reduction, producing an instance of size $cn$ for a constant $c$, and
solve it in $2^{o(cn)} = 2^{o(n)}$ time. The polynomial cost of the reduction
is absorbed.
\end{proof}

There is one obstacle left, and it is the reason the following lemma exists. A
$3$-SAT formula on $n$ variables may have as many as $\Theta(n^3)$ clauses, and
a reduction usually builds a gadget per clause, so the output has size
$\Theta(n^3)$ and is not linear in $n$. The fix is a remarkable theorem saying
that the clauses may as well be few.

\begin{theorem}[Sparsification lemma, \citet{ImpagliazzoPaturiZane2001}]
\label{thm:sparsification}
For every $\varepsilon > 0$ and every $k$ there is a constant $C =
C(k,\varepsilon)$ and an algorithm that, given a $k$-CNF formula $F$ on $n$
variables, outputs in $\bigoh{2^{\varepsilon n}\poly(n)}$ time formulas
$F_1, \ldots, F_t$ with $t \le 2^{\varepsilon n}$, each a $k$-CNF on the same
$n$ variables and with at most $Cn$ clauses, such that
\[ F \text{ is satisfiable} \iff F_i \text{ is satisfiable for some } i. \]
\end{theorem}

\begin{corollary}\label{cor:eth-sparse}
$\mathrm{ETH}$ implies that $3$-SAT has no algorithm running in
$2^{o(n+m)}$ time, where $m$ is the number of clauses.
\end{corollary}

\begin{proof}
Suppose such an algorithm exists and let $\varepsilon > 0$. Apply
\cref{thm:sparsification} and solve each of the at most $2^{\varepsilon n}$
sparse formulas, each of which has $n$ variables and at most $Cn$ clauses, in
$2^{o(n + Cn)} = 2^{o(n)}$ time. The total is
$2^{\varepsilon n} \cdot 2^{o(n)} \le 2^{2\varepsilon n}$ for large $n$. Since
$\varepsilon$ was arbitrary, $s_3 = 0$.
\end{proof}

\begin{intuition}
The sparsification lemma is doing something that sounds impossible: it removes
clauses without changing satisfiability. It does not, of course --- it splits
the formula into many formulas, and the point is that the number of pieces is
sub-exponential while each piece is sparse. Read it as a licence: \emph{when
proving lower bounds under ETH, you may assume $m = \bigoh{n}$}. Almost every
reduction in this chapter and the literature uses it in exactly that way, and
almost none uses anything else about it.
\end{intuition}

%==============================================================================
\section{No subexponential algorithm for 3-colouring}
\label{sec:eth-3col}
%==============================================================================

The first lower bound. \Cref{ch:exact-i} gave a $\bigoh{1.3217^n}$ algorithm
for $3$-colouring, the last in a long chain of improvements; the following says
the chain cannot converge to $1$.

\begin{theorem}\label{thm:eth-3col}
Assuming $\mathrm{ETH}$, there is no algorithm deciding \lang{3-Coloring} on a
graph with $n$ vertices in $2^{o(n)}$ time. The same holds with $n$ replaced by
the number of edges.
\end{theorem}

The proof is the classical reduction from $3$-SAT, checked for waste. We give
it in full, because the only new content of a fine-grained lower bound is the
counting, and the counting is invisible unless the construction is on the page.

\begin{proof}
Let $F$ be a $3$-CNF formula with variables $x_1,\ldots,x_n$ and clauses
$C_1, \ldots, C_m$. By \cref{cor:eth-sparse} we may assume $m = \bigoh{n}$.
Build a graph $G$ as follows.

\emph{The palette.} Take a triangle on three vertices named
$\mathsf{T}$, $\mathsf{F}$ and $\mathsf{B}$. In any $3$-colouring they receive
three different colours, and we name the colours after them: the colour of
$\mathsf{T}$ is ``true'', that of $\mathsf{F}$ is ``false'', that of
$\mathsf{B}$ is ``base''.

\emph{Variables.} For each variable $x_i$ add two vertices $v_i$ and
$\bar{v}_i$ and make $\{v_i, \bar v_i, \mathsf{B}\}$ a triangle. Then $v_i$ and
$\bar v_i$ get different colours, and neither is the base colour, so one of
them is true and the other false. Reading $v_i$ as $x_i$ and $\bar v_i$ as
$\neg x_i$, a colouring of these vertices is exactly a truth assignment.

\emph{Clauses.} For a clause $C = \ell_1 \vee \ell_2 \vee \ell_3$ add the
six-vertex gadget of \cref{fig:eth-clause-gadget}: two triangles
$\{a_1,a_2,a_3\}$ and $\{b_1,b_2,b_3\}$ wired so that $a_j$ is adjacent to the
vertex of $\ell_j$ for $j = 1,2$, that $a_1a_2b_1$ and $b_1b_2b_3$ are
triangles, that $b_2$ is adjacent to $\mathsf{F}$, and that $b_3$ is adjacent
to the vertex of $\ell_3$. The gadget has the property that it can be
$3$-coloured, given the colours of the three literal vertices, if and only if
at least one literal vertex has the true colour.

\emph{Correctness.} If $F$ has a satisfying assignment, colour the literal
vertices accordingly and extend over each clause gadget, which is possible
because each clause has a true literal. Conversely a $3$-colouring of $G$
assigns each variable a truth value by the colour of $v_i$, and no clause can
be all-false, so the assignment satisfies $F$.

\emph{The counting.} $G$ has $3 + 2n + 6m$ vertices and $\bigoh{n+m}$ edges. As
$m = \bigoh{n}$, both are $\bigoh{n}$: the reduction is linear. So an algorithm
deciding $3$-colouring in $2^{o(|V|)}$ time would decide $3$-SAT in $2^{o(n)}$
time, contradicting ETH by \cref{lem:linear-transfer}. The statement for edges
follows because $|E| = \bigoh{|V|}$ here.
\end{proof}

\begin{figure}[ht]
\centering
\figplaceholder{5.5cm}{The six-vertex clause gadget: two triangles, the three
literal vertices on the left attached to $a_1$, $a_2$ and $b_3$, the vertex
$\mathsf{F}$ of the palette attached to $b_2$. Beside it, the same gadget with
all three literals coloured false, and the resulting conflict marked.}
\caption{The clause gadget of \cref{thm:eth-3col}. It extends to a
$3$-colouring exactly when one of its three literal vertices is true; the
picture on the right shows the collision that occurs otherwise.}
\label{fig:eth-clause-gadget}
\end{figure}

\begin{remark}[What the bound does and does not say]\label{rem:eth-3col-scope}
\Cref{thm:eth-3col} rules out $2^{o(n)}$. It says nothing about the base: an
algorithm running in $1.0001^n$ time is not excluded by it, and no known
hypothesis excludes one. Fine-grained lower bounds of the form ``no
$\bigoh{(2-\varepsilon)^n}$'' do exist for some problems --- under SETH, not
ETH --- but $3$-colouring is not among them.

The bound does transfer, and cheaply. Any problem to which \lang{3-Coloring}
reduces linearly inherits it. \Cref{ch:planar} gave a $2^{\bigoh{\sqrt n}}$
algorithm for independent set on planar graphs, and the same technique gives one
for planar $3$-colouring; \cref{ex:eth-planar} asks you to show that ETH makes
the $\sqrt n$ optimal too.
\end{remark}

%==============================================================================
\section{Orthogonal vectors}
\label{sec:ov}
%==============================================================================

SETH is a statement about exponential time, and the lower bounds we now want
are about \emph{polynomial} time: we want to say that a problem with a
quadratic algorithm has no $\bigoh{n^{1.99}}$ one. The bridge between the two
is a single problem, which serves the fine-grained world the way $3$-SAT serves
\textrm{NP}-hardness.

\begin{probdef}{Orthogonal Vectors (OV)}
  \instance Two sets $A, B \subseteq \{0,1\}^d$ with $|A| = |B| = N$.
  \question Are there $a \in A$ and $b \in B$ with
  $\sum_{i=1}^{d} a_i b_i = 0$, that is, with no coordinate where both are
  $1$?
\end{probdef}

The trivial algorithm compares every pair, in $\bigoh{N^2 d}$ time. Nothing
substantially better is known when $d$ is, say, $\log^2 N$.

\begin{definition}\label{def:ov-conjecture}
The \emph{OV conjecture} states that for every $\delta > 0$ there is no
algorithm solving \lang{Orthogonal Vectors} in $\bigoh{N^{2-\delta}\poly(d)}$
time.
\end{definition}

\begin{theorem}[\citet{Williams2005}]\label{thm:seth-ov}
$\mathrm{SETH}$ implies the OV conjecture.
\end{theorem}

\begin{proof}
Let $\delta > 0$ and suppose \lang{Orthogonal Vectors} is solvable in
$\bigoh{N^{2-\delta}\poly(d)}$ time. Let $F$ be a CNF formula with $n$
variables and $m$ clauses; we solve it faster than $2^n$.

Split the variables into two halves $X_1$ and $X_2$ of $n/2$ each. For every
one of the $2^{n/2}$ assignments $\alpha$ to $X_1$, build a vector
$a_\alpha \in \{0,1\}^m$ with
\[ a_\alpha[j] = \begin{cases}
     0 & \text{if $\alpha$ already satisfies clause $C_j$,} \\
     1 & \text{otherwise;}
   \end{cases} \]
so a coordinate is $1$ exactly when that clause still needs help. Do the same
for the $2^{n/2}$ assignments $\beta$ to $X_2$, giving vectors $b_\beta$. Let
$A$ and $B$ be the two resulting sets, of size $N = 2^{n/2}$ in dimension
$d = m$.

Now $a_\alpha$ and $b_\beta$ are orthogonal if and only if there is no clause
that both halves fail to satisfy --- that is, if and only if the combined
assignment $\alpha \cup \beta$ satisfies $F$. So $F$ is satisfiable exactly
when $A$ and $B$ contain an orthogonal pair.

Building the vectors costs $\bigoh{2^{n/2} m}$ and the assumed algorithm then
costs
\[ \bigoh{\bigl(2^{n/2}\bigr)^{2-\delta} \poly(m)}
   = \bigoh{2^{(1 - \delta/2)n}\poly(m)}. \]
This is an algorithm for CNF-SAT with no bound on the clause width, running in
$\bigoh{2^{\delta' n}\poly(m)}$ for the fixed constant $\delta' = 1-\delta/2 <
1$. In particular it solves $k$-SAT within that bound for every $k$, so
$s_k \le \delta'$ for all $k$ and $\lim_k s_k \le \delta' < 1$, contradicting
SETH.
\end{proof}

\begin{intuition}
The proof is one idea: \emph{meet in the middle, and let the vector coordinates
be the clauses}. Splitting the variables in half is the standard trick of
\cref{ch:exact-i}, and it converts one exponential search into a quadratic
search over two exponentially large sets. Everything after that is
bookkeeping. What makes OV the right target is that this bookkeeping is
\emph{all} the structure that survives: an OV instance remembers only which
clauses each half-assignment leaves unsatisfied, and it turns out that this is
enough to encode a startling number of geometric and string problems.
\end{intuition}

%==============================================================================
\section{The Fr\'echet distance}
\label{sec:frechet}
%==============================================================================

Now the geometry. The Fr\'echet distance measures the similarity of two curves,
and it is the standard tool wherever trajectories are compared: matching GPS
tracks to a road network, comparing handwriting, clustering the paths of
migrating birds.

\begin{definition}[Discrete Fr\'echet distance]\label{def:discrete-frechet}
Let $P = (p_1,\ldots,p_n)$ and $Q = (q_1,\ldots,q_n)$ be sequences of points.
A \emph{traversal} is a sequence of pairs of indices starting at $(1,1)$ and
ending at $(n,n)$, in which each step advances the first index, the second
index, or both, by one. The \emph{discrete Fr\'echet distance}
$d_{dF}(P,Q)$ is the minimum over all traversals of the largest distance
$\lVert p_i - q_j \rVert$ occurring in it.
\end{definition}

The usual image is a person walking along $P$ and a dog along $Q$, joined by a
leash: both may move forward at any speed and neither may go back, and the
Fr\'echet distance is the shortest leash that suffices. The continuous version
lets them stand anywhere on the curve rather than only at vertices.

\begin{proposition}\label{prop:frechet-dp}
$d_{dF}(P,Q)$ can be computed in $\bigoh{n^2}$ time.
\end{proposition}

\begin{proof}
Fill an $n \times n$ table $D$, where $D[i][j]$ is the smallest leash length
that suffices for a traversal from $(1,1)$ to $(i,j)$:
\[ D[i][j] = \max\Bigl\{ \lVert p_i - q_j\rVert, \;
   \min\bigl( D[i-1][j],\, D[i][j-1],\, D[i-1][j-1] \bigr) \Bigr\}, \]
with $D[1][1] = \lVert p_1 - q_1 \rVert$. The answer is $D[n][n]$, and every
entry costs constant time.
\end{proof}

That algorithm is from 1994 and has never been improved beyond logarithmic
factors. The following says why.

\begin{theorem}[\citet{Bringmann2014}]\label{thm:frechet}
Assuming $\mathrm{SETH}$, there is no algorithm computing the discrete or the
continuous Fr\'echet distance of two curves with $n$ vertices in the plane in
$\bigoh{n^{2-\delta}}$ time, for any $\delta > 0$. The same holds for computing
a $1.001$-approximation of it.
\end{theorem}

We do not give the full construction --- it is a careful piece of gadgetry and
takes a paper --- but the shape of it is short, and the shape is the part worth
teaching.

\begin{proof}[Proof sketch]
The reduction is from \lang{Orthogonal Vectors}, so by
\cref{thm:seth-ov} it suffices to show that a strongly subquadratic Fr\'echet
algorithm gives a strongly subquadratic OV algorithm.

Given $A, B \subseteq \{0,1\}^d$ with $|A| = |B| = N$, one builds two curves,
$P$ from $A$ and $Q$ from $B$. Each vector $a \in A$ becomes a short piece
$P_a$ of the first curve, and each $b \in B$ a short piece $Q_b$ of the second,
each piece having $\bigoh{d}$ vertices; the pieces are strung together with
separating segments. The construction is arranged so that, for a fixed
threshold $\rho$,
\[ \text{the leash of length $\rho$ suffices to traverse $P_a$ against $Q_b$}
   \iff a \text{ and } b \text{ are orthogonal}, \]
and so that the separators force the traversal to line the pieces up one
against another rather than skipping. Then
\[ d_{dF}(P,Q) \le \rho \iff \text{some } a \in A, b \in B
   \text{ are orthogonal}. \]

The curves have $n = \bigoh{Nd}$ vertices. An
$\bigoh{n^{2-\delta}}$ algorithm for the Fr\'echet distance therefore decides
OV in $\bigoh{(Nd)^{2-\delta}} = \bigoh{N^{2-\delta}\poly(d)}$ time, which is
what \cref{thm:seth-ov} forbids.
\end{proof}

\begin{figure}[ht]
\centering
\figplaceholder{5cm}{Two polygonal curves in the plane with a traversal drawn
as a sequence of leash segments between them, beside the $n \times n$ table of
\cref{prop:frechet-dp} with a monotone staircase path through it.}
\caption{A traversal of two curves and the corresponding monotone path in the
dynamic programming table. The Fr\'echet distance is the minimum over such
paths of the largest leash used --- a bottleneck shortest path.}
\label{fig:frechet}
\end{figure}

\begin{intuition}
Notice what \cref{thm:frechet} costs and what it buys. It costs a hypothesis
that is stronger and less certain than $\mathrm{P} \ne \mathrm{NP}$. It buys a
statement about a problem with a perfectly good polynomial algorithm --- and
that is something no classical complexity theory can say at all. This is the
whole programme of fine-grained complexity: give up on unconditional bounds,
fix a small number of hypotheses, and organise the polynomial world by
reductions between concrete problems in the way the \textrm{NP}-hard world is
organised by reductions to \lang{SAT}.
\end{intuition}

%==============================================================================
\section{Notes and further reading}
\label{sec:lower-bounds-further-reading}
%==============================================================================

ETH and SETH are \citet{ImpagliazzoPaturi2001}; the sparsification lemma is
\citet{ImpagliazzoPaturiZane2001}. \Cref{thm:seth-ov} is due to
\citet{Williams2005}, where it appears as a step in an algorithm rather than as
a hardness result. \Cref{thm:frechet} is \citet{Bringmann2014}; the companion
result for edit distance, proved the same way, is
\citet{BackursIndyk2018}.

\citet[Chapter~14]{CyganEtAl2015} is the textbook account of ETH-based lower
bounds, including the planar and treewidth cases mentioned above.
\Cref{thm:tw-dominating-set} in \cref{ch:treewidth} is the treewidth example:
\citet{LokshtanovMarxSaurabh2018} show that SETH forbids an
$\bigoh{(3-\varepsilon)^k}$ algorithm for dominating set on a tree
decomposition of width $k$, which is why the three states of that dynamic
program are not an artefact. For clique,
\citet{ChenHuangKanjXia2006} show that ETH forbids an
$f(k)\, n^{o(k)}$ algorithm, so the $\bigoh{n^{k+1}}$ of
\cref{ex:fpt-three} is essentially optimal too.

%==============================================================================
\section{Exercises}
\label{sec:ex-lower-bounds}
%==============================================================================

\begin{exercise}
Show that ETH implies \lang{Vertex Cover} has no $2^{o(n)}$ algorithm, where
$n$ is the number of vertices. Which reduction do you use, and is it linear?
\end{exercise}

\begin{exercise}
The standard reduction from $3$-SAT to \lang{Independent Set} builds a triangle
per clause and joins contradictory literals. Check that it is linear, and state
the consequence.
\end{exercise}

\begin{exercise}\label{ex:eth-planar}
\Cref{ch:planar} solves independent set on planar graphs in
$2^{\bigoh{\sqrt n}}$ time. Show that, assuming ETH, there is no
$2^{o(\sqrt n)}$ algorithm. Hint: the standard reduction to the planar case
replaces each crossing by a constant-size gadget, and a graph with $\bigoh{n}$
edges can be drawn with $\bigoh{n^2}$ crossings.
\end{exercise}

\begin{exercise}
Show that ETH implies that $3$-SAT has no $2^{o(m)}$ algorithm, where $m$ is
the number of clauses. Where is the sparsification lemma used?
\end{exercise}

\begin{exercise}
In the proof of \cref{thm:seth-ov} the two sets have equal size because the
variables were split evenly. What running time does the same argument give if
the split is $\gamma n$ and $(1-\gamma)n$, and why is $\gamma = 1/2$ the right
choice?
\end{exercise}

\begin{exercise}
\lang{Longest Common Subsequence} of two strings of length $n$ has a textbook
$\bigoh{n^2}$ algorithm. Assuming the OV conjecture, would you expect an
$\bigoh{n^{1.9}}$ one? Look up the answer and identify which step of the
argument of \cref{thm:frechet} it replaces.
\end{exercise}

\begin{exercise}
Suppose someone announces a $\bigoh{1.99^n}$ algorithm for CNF-SAT. Which
statements in this chapter survive, and which are lost?
\end{exercise}

\begin{exercise}
The dynamic program of \cref{prop:frechet-dp} computes a bottleneck shortest
path in a grid. Give an $\bigoh{n^2 \log n}$ algorithm for the same problem by
binary search on the answer, and explain why this does not contradict
\cref{thm:frechet}.
\end{exercise}
